% =====================================================================
%  Banco complementar --- Lei de Coulomb  (REVISADO)
%  Substitui a versao gerada por template: 32 questoes, todas com
%  enunciado distinto. Nivel 4 reclassificado conforme o criterio da
%  obra (sintese, demonstracao, otimizacao ou modelagem nao rotineira).
%  Distribuicao mantida: 7 N1 + 10 N2 + 6 N3 + 9 N4 = 32
% =====================================================================

\subsubsection*{Banco complementar --- Lei de Coulomb}

\begin{atencao}[Organização do banco]
As questões abaixo complementam o tópico sem repetir enunciados da seção
principal. Cada nível explora uma \emph{estrutura de raciocínio} diferente:
N1 aplica a fórmula em suas quatro formas isoladas ($F$, $q$, $r$, meio);
N2 exige composição ou reinterpretação; N3 combina vetores, equilíbrio e
redistribuição de carga; N4 pede demonstração, integração, otimização ou
modelagem.
\end{atencao}

\blocofixacao

\begin{questao}[1]{}
Duas cargas puntiformes $q_1=\SI{20}{\nano\coulomb}$ e
$q_2=\SI{30}{\nano\coulomb}$ estão separadas por $\SI{3}{\centi\meter}$ no
vácuo. Determine o módulo da força elétrica e classifique-a.
\dados{$\ke=\SI{9e9}{\newton\meter\squared\per\coulomb\squared}$}
\resposta{$F=\SI{6}{\milli\newton}$, repulsiva}
\resolucao{Convertendo antes de substituir --- este é o erro mais comum:
$q_1=\pot{2}{-8}\,\si{\coulomb}$, $q_2=\pot{3}{-8}\,\si{\coulomb}$,
$r=\pot{3}{-2}\,\si{\meter}$.
\[
F=\ke\frac{q_1q_2}{r^{2}}
 =\frac{\pot{9}{9}\cdot\pot{2}{-8}\cdot\pot{3}{-8}}{(\pot{3}{-2})^{2}}
 =\frac{\pot{5.4}{-6}}{\pot{9}{-4}}=\pot{6}{-3}\,\si{\newton}.
\]
Ambas positivas $\Rightarrow$ repulsão.}
\end{questao}

\begin{questao}[1]{}
Duas esferas idênticas recebem cargas iguais $q$ e se repelem com força de
$\SI{0.90}{\newton}$ quando separadas por $\SI{20}{\centi\meter}$.
\begin{itensq}
\item Determine $q$.
\item Quantos elétrons foram \textbf{retirados} de cada esfera?
\end{itensq}
\dados{$e=\pot{1.6}{-19}\,\si{\coulomb}$}
\resposta{(a) $q=\SI{2}{\micro\coulomb}$; (b) $n=\pot{1.25}{13}$ elétrons}
\resolucao{(a) De $F=\ke q^{2}/r^{2}$,
\[
q=\sqrt{\frac{Fr^{2}}{\ke}}=\sqrt{\frac{0{,}90\cdot(0{,}20)^{2}}{\pot{9}{9}}}
 =\sqrt{\pot{4}{-12}}=\pot{2}{-6}\,\si{\coulomb}.
\]
(b) Repulsão com cargas iguais e sinal positivo indica \emph{falta} de
elétrons: $n=q/e=\pot{2}{-6}/\pot{1.6}{-19}=\pot{1.25}{13}$.}
\end{questao}

\begin{questao}[1]{}
Duas cargas de $\SI{5}{\micro\coulomb}$ cada devem exercer uma sobre a outra
força de $\SI{2.5}{\newton}$ no vácuo. A que distância devem ser colocadas?
\resposta{$r=\SI{0.30}{\meter}$}
\resolucao{Isolando $r$ na lei de Coulomb:
\[
r=\sqrt{\frac{\ke q_1q_2}{F}}
 =\sqrt{\frac{\pot{9}{9}\cdot(\pot{5}{-6})^{2}}{2{,}5}}
 =\sqrt{\frac{0{,}225}{2{,}5}}=\sqrt{0{,}09}=\SI{0.30}{\meter}.
\]}
\end{questao}

\begin{questao}[1]{}
Duas cargas puntiformes se atraem com força $F$. Mantendo as cargas e
\textbf{dobrando} a distância entre elas, a nova força será:
\begin{alternativas}[3]
\item $F/4$
\item $F/2$
\item $F$
\item $2F$
\item $4F$
\end{alternativas}
\resposta{(a)}
\resolucao{$F\propto 1/r^{2}$. Dobrar $r$ divide a força por $2^{2}=4$:
$F'=F/4$. Note que a natureza (atrativa) não muda --- ela depende apenas dos
sinais, nunca da distância.}
\end{questao}

\begin{questao}[1]{}
Duas cargas separadas por uma distância fixa exercem força de
$\SI{0.80}{\newton}$ no vácuo. Todo o conjunto é então imerso em um líquido
isolante de permissividade relativa $\varepsilon_r=4$. Determine a nova força.
\resposta{$F'=\SI{0.20}{\newton}$}
\resolucao{Em um meio dielétrico, $F_{\text{meio}}=F_{\text{vácuo}}/\varepsilon_r$,
pois $\ke$ é substituída por $\ke/\varepsilon_r$:
\[
F'=\frac{0{,}80}{4}=\SI{0.20}{\newton}.
\]
O dielétrico se polariza e blinda parcialmente as cargas --- por isso a força
\emph{sempre} diminui ($\varepsilon_r\ge 1$).}
\end{questao}

\begin{questao}[1]{}
Um bastão carregado positivamente é aproximado --- sem tocar --- de uma esfera
metálica \textbf{neutra} e isolada. A força entre eles é:
\begin{alternativas}[2]
\item nula, pois a esfera é neutra
\item atrativa
\item repulsiva
\item atrativa ou repulsiva, dependendo da distância
\end{alternativas}
\resposta{(b)}
\resolucao{A carga \emph{líquida} da esfera continua nula, mas o bastão
positivo atrai elétrons livres para a face próxima e repele carga positiva
para a face distante (\textbf{indução eletrostática}). A face próxima, negativa,
está \emph{mais perto} do bastão que a face distante, positiva. Como
$F\propto 1/r^{2}$, a atração vence: resultante \textbf{atrativa}. É o mesmo
princípio pelo qual um canudo atritado atrai pedaços de papel neutros.}
\end{questao}

\begin{questao}[1]{}
Estime a força entre duas cargas de $\SI{1}{\coulomb}$ separadas por
$\SI{1}{\kilo\meter}$ e comente o resultado.
\resposta{$F=\SI{9}{\kilo\newton}$ (o peso de cerca de \SI{900}{\kilogram})}
\resolucao{
\[
F=\frac{\pot{9}{9}\cdot 1\cdot 1}{(10^{3})^{2}}=\pot{9}{3}\,\si{\newton}.
\]
Nove mil newtons --- o peso de quase uma tonelada --- a \emph{um quilômetro}
de distância. Isso mostra por que o coulomb é uma unidade enorme para a
eletrostática: cargas realmente separáveis por atrito ficam na faixa de
$\si{\nano\coulomb}$ a $\si{\micro\coulomb}$.}
\end{questao}

\blocoaplicacao

\begin{questao}[2]{}
Três cargas estão fixas sobre o eixo $x$: $q_1=\SI{+3}{\micro\coulomb}$ em
$x=0$, $q_2=\SI{+2}{\micro\coulomb}$ em $x=\SI{20}{\centi\meter}$ e
$q_3=\SI{-5}{\micro\coulomb}$ em $x=\SI{50}{\centi\meter}$. Determine a força
resultante sobre $q_2$ (módulo e sentido).
\resposta{$F=\SI{2.35}{\newton}$, no sentido $+x$}
\resolucao{Trate cada par separadamente e some \emph{com sinal}.
\[
F_{12}=\frac{\pot{9}{9}\cdot\pot{3}{-6}\cdot\pot{2}{-6}}{(0{,}20)^{2}}
=\SI{1.35}{\newton}\quad(\text{repulsão}\Rightarrow +x)
\]
\[
F_{32}=\frac{\pot{9}{9}\cdot\pot{2}{-6}\cdot\pot{5}{-6}}{(0{,}30)^{2}}
=\SI{1.00}{\newton}\quad(\text{atração por } q_3 \Rightarrow +x)
\]
As duas apontam no mesmo sentido: $F=1{,}35+1{,}00=\SI{2.35}{\newton}$ em $+x$.
\textbf{Cuidado:} não some as cargas antes; a lei de Coulomb vale par a par
(princípio da superposição).}
\end{questao}

\begin{questao}[2]{}
Duas esferas metálicas idênticas com $q_1=\SI{+8}{\micro\coulomb}$ e
$q_2=\SI{-2}{\micro\coulomb}$ estão a $\SI{30}{\centi\meter}$. Elas são postas
em contato e recolocadas na mesma posição.
\begin{itensq}
\item Calcule a força antes do contato.
\item Calcule a força depois.
\end{itensq}
\resposta{(a) $\SI{1.6}{\newton}$, atrativa; (b) $\SI{0.9}{\newton}$, repulsiva}
\resolucao{(a) $F=\dfrac{\pot{9}{9}\cdot\pot{8}{-6}\cdot\pot{2}{-6}}{(0{,}30)^{2}}
=\dfrac{0{,}144}{0{,}09}=\SI{1.6}{\newton}$, atrativa (sinais opostos).\\
(b) Esferas \emph{idênticas} em contato dividem a carga total igualmente:
\[
q'=\frac{q_1+q_2}{2}=\frac{+8-2}{2}=\SI{+3}{\micro\coulomb}\ \text{cada}.
\]
\[
F'=\frac{\pot{9}{9}\cdot(\pot{3}{-6})^{2}}{0{,}09}=\SI{0.9}{\newton},
\ \text{agora repulsiva}.
\]
A força \emph{e} a natureza mudaram: o contato conserva a carga total, não a
individual.}
\end{questao}

\begin{questao}[2]{}
Duas cargas fixas exercem entre si $\SI{0.60}{\newton}$ no vácuo. O conjunto é
mergulhado em óleo com $\varepsilon_r=2{,}5$.
\begin{itensq}
\item Qual é a nova força, mantida a distância?
\item Que fração da distância original restauraria a força de
      $\SI{0.60}{\newton}$?
\end{itensq}
\resposta{(a) $\SI{0.24}{\newton}$; (b) $r'=r/\sqrt{2{,}5}\approx 0{,}63\,r$}
\resolucao{(a) $F'=F/\varepsilon_r=0{,}60/2{,}5=\SI{0.24}{\newton}$.\\
(b) Queremos $\dfrac{\ke q_1q_2}{\varepsilon_r r'^{2}}=\dfrac{\ke q_1q_2}{r^{2}}$,
ou seja $\varepsilon_r r'^{2}=r^{2}$:
\[
r'=\frac{r}{\sqrt{\varepsilon_r}}=\frac{r}{\sqrt{2{,}5}}=0{,}632\,r.
\]
Aproximar as cargas em cerca de $37\%$ compensa exatamente o efeito do meio.}
\end{questao}

\begin{questao}[2]{}
A força entre duas cargas $q_1$ e $q_2$ separadas por $d$ vale $F$. Se $q_1$ for
\textbf{triplicada} e a distância for reduzida à \textbf{metade}, a nova força
será:
\begin{alternativas}[3]
\item $1{,}5F$
\item $3F$
\item $6F$
\item $12F$
\item $24F$
\end{alternativas}
\resposta{(d)}
\resolucao{Os efeitos são multiplicativos e independentes:
\[
F'=\ke\frac{(3q_1)q_2}{(d/2)^{2}}=3\cdot 4\cdot \ke\frac{q_1q_2}{d^{2}}=12F.
\]
Fator $3$ da carga $\times$ fator $4$ da distância $=12$.}
\end{questao}

\begin{questao}[2]{}
Em uma impressora a jato eletrostático, cada gota de tinta tem massa
$\pot{1.0}{-11}\,\si{\kilogram}$ e é carregada pela remoção de
$\pot{2.0}{5}$ elétrons. Duas gotas vizinhas ficam a $\SI{1.0}{\milli\meter}$.
Compare a força elétrica entre elas com o peso de uma gota.
\dados{$e=\pot{1.6}{-19}\,\si{\coulomb}$; $g=\SI{9.8}{\meter\per\second\squared}$}
\resposta{$F_e=\pot{9.2}{-12}\,\si{\newton}$; $P=\pot{9.8}{-11}\,\si{\newton}$;
$F_e/P\approx 0{,}094$}
\resolucao{Carga da gota: $q=ne=\pot{2.0}{5}\cdot\pot{1.6}{-19}
=\pot{3.2}{-14}\,\si{\coulomb}$.
\[
F_e=\frac{\pot{9}{9}\cdot(\pot{3.2}{-14})^{2}}{(\pot{1.0}{-3})^{2}}
=\pot{9.2}{-12}\,\si{\newton}.
\]
Peso: $P=mg=\pot{1.0}{-11}\cdot 9{,}8=\pot{9.8}{-11}\,\si{\newton}$.
Razão $\approx 9{,}4\%$: a repulsão entre gotas é pequena diante do peso, mas
não desprezível --- é ela que degrada a nitidez quando as gotas são emitidas
muito próximas.}
\end{questao}

\begin{questao}[2]{}
Duas cargas positivas, $q_1=\SI{4}{\micro\coulomb}$ em $x=0$ e
$q_2=\SI{9}{\micro\coulomb}$ em $x=\SI{50}{\centi\meter}$, estão fixas.
Determine o ponto onde uma terceira carga ficaria em equilíbrio.
\resposta{$x=\SI{20}{\centi\meter}$}
\resolucao{Como as duas são positivas, o ponto de força nula está
\emph{entre} elas. Seja $x$ a distância a $q_1$:
\[
\frac{\ke q_1 q_3}{x^{2}}=\frac{\ke q_2 q_3}{(0{,}50-x)^{2}}
\;\Longrightarrow\;
\frac{4}{x^{2}}=\frac{9}{(0{,}50-x)^{2}}.
\]
Extraindo a raiz (ambos os lados positivos):
$\dfrac{2}{x}=\dfrac{3}{0{,}50-x}\Rightarrow 2(0{,}50-x)=3x
\Rightarrow x=\SI{0.20}{\meter}$.
O ponto fica \emph{mais próximo da carga menor}, como esperado. Note que $q_3$
se cancela: a posição não depende dela.}
\end{questao}

\begin{questao}[2]{}
Uma carga $q$ e outra $100q$ interagem. Sobre os módulos das forças que uma
exerce sobre a outra, é correto afirmar:
\begin{alternativas}[2]
\item a força sobre $q$ é 100 vezes maior
\item a força sobre $100q$ é 100 vezes maior
\item as forças têm o mesmo módulo
\item depende das massas
\end{alternativas}
\resposta{(c)}
\resolucao{A lei de Coulomb é simétrica no produto $q_1q_2$: ambas valem
$\ke\,q\cdot 100q/r^{2}$. Formam um par ação--reação (3\textsuperscript{a} lei
de Newton): mesmo módulo, mesma direção, sentidos opostos. O que \emph{difere}
é a \textbf{aceleração}, se as massas forem diferentes --- confundir força com
aceleração é o erro clássico aqui.}
\end{questao}

\begin{questao}[2]{}
No modelo de Bohr para o hidrogênio, o elétron orbita o próton a
$\pot{5.3}{-11}\,\si{\meter}$. Determine a força elétrica e a aceleração
centrípeta do elétron.
\dados{$e=\pot{1.6}{-19}\,\si{\coulomb}$; $m_e=\pot{9.11}{-31}\,\si{\kilogram}$}
\resposta{$F=\pot{8.2}{-8}\,\si{\newton}$;
$a=\pot{9.0}{22}\,\si{\meter\per\second\squared}$}
\resolucao{
\[
F=\frac{\pot{9}{9}\cdot(\pot{1.6}{-19})^{2}}{(\pot{5.3}{-11})^{2}}
=\frac{\pot{2.304}{-28}}{\pot{2.81}{-21}}=\pot{8.2}{-8}\,\si{\newton}.
\]
\[
a=\frac{F}{m_e}=\frac{\pot{8.2}{-8}}{\pot{9.11}{-31}}
=\pot{9.0}{22}\,\si{\meter\per\second\squared}.
\]
Uma força minúscula produz aceleração de $10^{22}\,\si{\meter\per\second\squared}$
porque a massa do elétron é ainda mais minúscula.}
\end{questao}

\begin{questao}[2]{}
Duas esferas metálicas idênticas com cargas $+Q$ e $+3Q$, separadas por $d$,
repelem-se com força $F$. Elas são postas em contato e depois afastadas até
$2d$. Determine a nova força em função de $F$.
\resposta{$F'=F/3$}
\resolucao{Antes: $F=\ke\dfrac{Q\cdot 3Q}{d^{2}}=\dfrac{3\ke Q^{2}}{d^{2}}$.\\
Após o contato, cada esfera fica com $\dfrac{Q+3Q}{2}=2Q$. A $2d$:
\[
F'=\ke\frac{(2Q)(2Q)}{(2d)^{2}}=\ke\frac{4Q^{2}}{4d^{2}}=\frac{\ke Q^{2}}{d^{2}}
=\frac{F}{3}.
\]
O produto das cargas subiu (de $3Q^2$ para $4Q^2$), mas o fator $4$ da
distância dominou.}
\end{questao}

\begin{questao}[2]{}
Compare a força entre duas cargas de $\SI{1.0}{\micro\coulomb}$ separadas por
$\SI{1.0}{\centi\meter}$ com o peso de um corpo de $\SI{1.0}{\kilogram}$, e
interprete.
\dados{$g=\SI{9.8}{\meter\per\second\squared}$}
\resposta{$F=\SI{90}{\newton}$ contra $P=\SI{9.8}{\newton}$; razão $\approx 9{,}2$}
\resolucao{
\[
F=\frac{\pot{9}{9}\cdot(\pot{1}{-6})^{2}}{(\pot{1}{-2})^{2}}
=\frac{\pot{9}{-3}}{\pot{1}{-4}}=\SI{90}{\newton};
\qquad P=1{,}0\cdot 9{,}8=\SI{9.8}{\newton}.
\]
Um milionésimo de coulomb, a um centímetro, produz nove vezes o peso de um
quilograma. Essa desproporção entre a força elétrica e a gravitacional é a razão
pela qual é a matéria macroscópica é praticamente neutra: qualquer desequilíbrio
apreciável de carga seria violentamente desfeito.}
\end{questao}

\blocoaprofundamento

\begin{questao}[3]{}
Duas cargas fixas sobre o eixo $x$: $q_1=\SI{+4}{\micro\coulomb}$ em $x=0$ e
$q_2=\SI{-1}{\micro\coulomb}$ em $x=\SI{30}{\centi\meter}$. Determine todos os
pontos do eixo onde a força sobre uma carga de prova seria nula.
\resposta{Um único ponto: $x=\SI{60}{\centi\meter}$}
\resolucao{Com \textbf{sinais opostos}, o ponto não pode estar entre as cargas
(ali as duas forças apontam para o mesmo lado). Deve estar \emph{fora} do
segmento e, mais especificamente, do lado da carga de \emph{menor} módulo ---
só assim a maior proximidade compensa o menor valor. Testando $x>0{,}30$:
\[
\frac{\ke\cdot 4\mu}{x^{2}}=\frac{\ke\cdot 1\mu}{(x-0{,}30)^{2}}
\;\Longrightarrow\;
\frac{2}{x}=\frac{1}{x-0{,}30}
\;\Longrightarrow\;
2x-0{,}60=x \Rightarrow x=\SI{0.60}{\meter}.
\]
\textbf{Verificação:}
$\ke\cdot4\mu/(0{,}60)^{2}=\ke\cdot1\mu/(0{,}30)^{2}$ --- confere. A raiz
$x=0{,}20$ da equação quadrática completa é \emph{espúria}: corresponde ao
ponto interno, onde as forças se somam em vez de se cancelar.}
\end{questao}

\begin{questao}[3]{}
Seis cargas iguais $Q=\SI{2}{\micro\coulomb}$ ocupam os vértices de um hexágono
regular de lado $a=\SI{10}{\centi\meter}$, com uma carga
$q=\SI{1}{\micro\coulomb}$ no centro.
\begin{itensq}
\item Qual é a força resultante sobre $q$?
\item Uma das seis cargas é removida. Qual passa a ser a força sobre $q$
      (módulo e direção)?
\end{itensq}
\resposta{(a) nula; (b) $\SI{1.8}{\newton}$, apontando \emph{para} o vértice vazio}
\resolucao{(a) Em um hexágono regular, o centro dista $a$ de todos os vértices
e as cargas ocupam posições diametralmente opostas duas a duas. Cada par
contribui com forças de mesmo módulo e sentidos opostos: resultante nula, por
simetria.\\
(b) Use o argumento da superposição \emph{ao contrário}: 
$\vec{F}_{6}= \vec{F}_{5}+\vec{F}_{\text{removida}}=\vec{0}$, logo
$\vec{F}_{5}=-\vec{F}_{\text{removida}}$.
\[
|\vec{F}_{5}|=\frac{\ke Qq}{a^{2}}
=\frac{\pot{9}{9}\cdot\pot{2}{-6}\cdot\pot{1}{-6}}{(0{,}10)^{2}}
=\SI{1.8}{\newton}.
\]
A carga removida \emph{repelia} $q$ para longe do vértice; sem ela, a resultante
aponta \textbf{para} o vértice vazio. Este truque de simetria evita somar cinco
vetores.}
\end{questao}

\begin{questao}[3]{}
Três cargas ocupam os vértices de um triângulo retângulo:
$q_A=\SI{+2}{\micro\coulomb}$ na origem, $q_B=\SI{+3}{\micro\coulomb}$ em
$(\SI{30}{\centi\meter},0)$ e $q_C=\SI{+4}{\micro\coulomb}$ em
$(0,\SI{40}{\centi\meter})$. Determine a força resultante sobre $q_A$: módulo e
ângulo com o eixo $x$.
\resposta{$F=\SI{0.75}{\newton}$, a $\SI{216.9}{\degree}$ (isto é, $\SI{36.9}{\degree}$ abaixo
do semieixo $-x$)}
\resolucao{As duas forças sobre $q_A$ são perpendiculares --- o que torna a
composição imediata.
\[
F_{AB}=\frac{\pot{9}{9}\cdot\pot{2}{-6}\cdot\pot{3}{-6}}{(0{,}30)^{2}}
=\SI{0.60}{\newton}\ \ (\text{repulsão}\Rightarrow -x)
\]
\[
F_{AC}=\frac{\pot{9}{9}\cdot\pot{2}{-6}\cdot\pot{4}{-6}}{(0{,}40)^{2}}
=\SI{0.45}{\newton}\ \ (\text{repulsão}\Rightarrow -y)
\]
\[
F=\sqrt{0{,}60^{2}+0{,}45^{2}}=\sqrt{0{,}5625}=\SI{0.75}{\newton},
\qquad
\tan\theta=\frac{0{,}45}{0{,}60}=0{,}75 \Rightarrow \theta=\SI{36.87}{\degree}.
\]
Um clássico triângulo $3$--$4$--$5$ escondido. \textbf{Atenção:} a hipotenusa
$BC=\SI{50}{\centi\meter}$ é irrelevante --- ela governa a força \emph{entre}
$q_B$ e $q_C$, não a força sobre $q_A$.}
\end{questao}

\begin{questao}[3]{}
Duas pequenas esferas metálicas idênticas, separadas por $\SI{10}{\centi\meter}$,
\textbf{atraem-se} com força de $\SI{0.108}{\newton}$. Postas em contato e
recolocadas na mesma distância, passam a \textbf{repelir-se} com
$\SI{0.036}{\newton}$. Determine as cargas iniciais.
\resposta{$q_1=\SI{+0.6}{\micro\coulomb}$ e $q_2=\SI{-0.2}{\micro\coulomb}$
(ou a permuta)}
\resolucao{\textbf{Etapa 1 --- produto.} A atração indica sinais opostos:
\[
|q_1q_2|=\frac{F_1r^{2}}{\ke}=\frac{0{,}108\cdot 0{,}01}{\pot{9}{9}}
=\pot{1.2}{-13}\ \Rightarrow\ q_1q_2=-\pot{1.2}{-13}\,\si{\coulomb\squared}.
\]
\textbf{Etapa 2 --- soma.} Após o contato cada esfera fica com
$q'=(q_1+q_2)/2$:
\[
q'=\sqrt{\frac{F_2r^{2}}{\ke}}=\sqrt{\frac{0{,}036\cdot0{,}01}{\pot{9}{9}}}
=\pot{2}{-7}\,\si{\coulomb}
\ \Rightarrow\ q_1+q_2=\pot{4}{-7}\,\si{\coulomb}.
\]
\textbf{Etapa 3 --- soma e produto.} As cargas são raízes de
$x^{2}-Sx+P=0$ com $S=\pot{4}{-7}$ e $P=-\pot{1.2}{-13}$:
\[
x=\frac{\pot{4}{-7}\pm\sqrt{\pot{1.6}{-13}+\pot{4.8}{-13}}}{2}
 =\frac{\pot{4}{-7}\pm\pot{8}{-7}}{2}
 \Rightarrow x=\pot{6}{-7}\ \text{ou}\ -\pot{2}{-7}.
\]
A repulsão final confirma que a soma é não nula. Se as cargas fossem simétricas
($q$ e $-q$), o contato as neutralizaria e a força final seria zero.}
\end{questao}

\begin{questao}[3]{}
Um pequeno bloco de massa $m=\SI{20}{\gram}$, com carga $q=\SI{1}{\micro\coulomb}$,
repousa sobre um plano inclinado \textbf{liso} de $\SI{30}{\degree}$. Uma carga fixa
$Q=\SI{1}{\micro\coulomb}$ está presa na base do plano, e a força elétrica atua
ao longo da superfície. Determine a distância de equilíbrio $d$ medida sobre o
plano.
\dados{$g=\SI{9.8}{\meter\per\second\squared}$}
\resposta{$d\approx\SI{30.3}{\centi\meter}$}
\resolucao{Sobre o plano liso, o equilíbrio na direção da superfície envolve
apenas a componente do peso e a repulsão elétrica (a normal é perpendicular):
\[
\frac{\ke Qq}{d^{2}}=mg\sin\theta.
\]
\[
mg\sin\SI{30}{\degree}=0{,}020\cdot 9{,}8\cdot 0{,}5=\SI{0.098}{\newton};
\qquad
d=\sqrt{\frac{\pot{9}{9}\cdot(\pot{1}{-6})^{2}}{0{,}098}}
=\sqrt{0{,}0918}=\SI{0.303}{\meter}.
\]
\textbf{Estabilidade:} se o bloco descer, $d$ diminui e a repulsão cresce como
$1/d^{2}$, empurrando-o de volta; se subir, a repulsão cai e o peso o traz de
volta. O equilíbrio é \emph{estável}.}
\end{questao}

\begin{questao}[3]{}
Um dipolo é formado por $+q$ e $-q$, com $q=\SI{1}{\micro\coulomb}$, fixos em
$(\pm\SI{5}{\centi\meter},0)$. Uma carga $Q=\SI{2}{\micro\coulomb}$ é colocada
em $(0,\SI{12}{\centi\meter})$, sobre a mediatriz.
\begin{itensq}
\item Mostre que a resultante sobre $Q$ é paralela ao eixo do dipolo.
\item Calcule seu módulo.
\end{itensq}
\resposta{(a) as componentes verticais se cancelam; (b) $F=\SI{0.82}{\newton}$,
no sentido $-x$ (de $+q$ para $-q$)}
\resolucao{(a) Cada carga dista
$r=\sqrt{0{,}05^{2}+0{,}12^{2}}=\SI{0.13}{\meter}$ de $Q$ (triângulo
$5$--$12$--$13$). A carga $+q$ \emph{repele} $Q$ e $-q$ a \emph{atrai}: as duas
forças têm o mesmo módulo e as componentes em $y$ apontam em sentidos opostos,
cancelando-se. Restam as componentes em $x$, que se \textbf{somam}.\\
(b) Cada força vale $F_1=\ke qQ/r^{2}$ e sua componente horizontal é
$F_1\cdot(a/r)$, com $a=\SI{0.05}{\meter}$:
\[
F=2\,\frac{\ke qQ}{r^{2}}\cdot\frac{a}{r}=\frac{2\ke qQ\,a}{r^{3}}
=\frac{2\cdot\pot{9}{9}\cdot\pot{1}{-6}\cdot\pot{2}{-6}\cdot 0{,}05}
       {(0{,}13)^{3}}=\SI{0.82}{\newton}.
\]
Note a dependência $1/r^{3}$: longe do dipolo a força cai \emph{mais rápido}
que a de uma carga isolada, porque as duas cargas quase se cancelam.}
\end{questao}

\blocodesafio

\begin{questao}[4]{}
\textbf{(Modelagem experimental.)} Em uma reprodução da balança de torção, duas
esferas com cargas iguais $q=\SI{0.10}{\micro\coulomb}$ foram mantidas a várias
distâncias e a força foi medida:

\begin{center}
\begin{tabular}{cccccc}
\hline
$r$ (m) & 0,020 & 0,030 & 0,040 & 0,050 & 0,060\\
$F$ (N) & 0,223 & 0,101 & 0,0558 & 0,0362 & 0,0249\\
\hline
\end{tabular}
\end{center}

\begin{itensq}
\item Proponha uma linearização que permita testar a hipótese $F=C\,r^{n}$.
\item Estime $n$ e discuta se os dados sustentam a lei do inverso do quadrado.
\item Obtenha o valor experimental de $\ke$ e compare com o tabelado.
\end{itensq}
\resposta{(a) $\log F=\log C+n\log r$; (b) $n\approx-2{,}00$; (c)
$\ke\approx\pot{9.1}{9}\,\si{\newton\meter\squared\per\coulomb\squared}$
(erro $\approx 1\%$)}
\resolucao{(a) Tomando logaritmos, uma lei de potência vira uma reta:
$\log F = \log C + n\log r$. O coeficiente angular do gráfico
$\log F \times \log r$ é o expoente procurado.\\
(b) Ajustando os cinco pontos por mínimos quadrados obtém-se
$n=-1{,}997$ e $C=\pot{9.09}{-5}$. O expoente reproduz $-2$ com desvio de
$0{,}15\%$ --- dentro do que se espera de medidas de força dessa ordem. Um teste
rápido sem regressão: dobrar $r$ de $0{,}020$ para $0{,}040$ deveria dividir $F$
por $4$; de fato $0{,}223/0{,}0558=4{,}00$.\\
(c) Como $C=\ke q^{2}$ e $q^{2}=\pot{1}{-14}\,\si{\coulomb\squared}$:
\[
\ke=\frac{C}{q^{2}}=\frac{\pot{9.09}{-5}}{\pot{1}{-14}}
=\pot{9.09}{9}\,\si{\newton\meter\squared\per\coulomb\squared},
\]
contra $\pot{8.99}{9}$ tabelado --- erro relativo de $1{,}1\%$.\\
\textbf{Observação de método:} a linearização logarítmica é preferível a ajustar
$F\times 1/r^{2}$ diretamente, porque \emph{não pressupõe} o expoente --- ela o
mede. Esse é exatamente o argumento que Coulomb precisou construir em 1785.}
\end{questao}

\begin{questao}[4]{}
\textbf{(Otimização com cálculo.)} Um anel fino de raio $R$, com carga total $Q$
uniformemente distribuída, está no plano $yz$. Uma carga puntiforme $q$ é
colocada sobre o eixo do anel, a uma distância $x$ do centro.
\begin{itensq}
\item Mostre que $F(x)=\dfrac{\ke Q q\,x}{(x^{2}+R^{2})^{3/2}}$.
\item Determine o valor de $x$ que \textbf{maximiza} a força.
\item Calcule $F_{\max}$ para $Q=\SI{10}{\micro\coulomb}$,
      $q=\SI{2}{\micro\coulomb}$ e $R=\SI{10}{\centi\meter}$.
\end{itensq}
\resposta{(b) $x=R/\sqrt{2}\approx\SI{7.07}{\centi\meter}$;
(c) $F_{\max}=\dfrac{2\ke Qq}{3\sqrt{3}\,R^{2}}=\SI{6.93}{\newton}$}
\resolucao{(a) Divida o anel em elementos $\mathrm{d} Q$. Cada um dista
$s=\sqrt{x^{2}+R^{2}}$ de $q$ e contribui com $\ke q\,\mathrm{d} Q/s^{2}$. Por
simetria, as componentes perpendiculares ao eixo se cancelam aos pares; sobra a
projeção axial, fator $\cos\alpha=x/s$:
\[
F=\int \frac{\ke q\,\mathrm{d} Q}{s^{2}}\cdot\frac{x}{s}
 =\frac{\ke q x}{(x^{2}+R^{2})^{3/2}}\int \mathrm{d} Q
 =\frac{\ke Q q\,x}{(x^{2}+R^{2})^{3/2}}.
\]
(b) Derivando e igualando a zero:
\[
\frac{\mathrm{d} F}{\mathrm{d} x}\propto (x^{2}+R^{2})^{3/2}-x\cdot 3x(x^{2}+R^{2})^{1/2}=0
\;\Rightarrow\; x^{2}+R^{2}-3x^{2}=0 \;\Rightarrow\; x=\frac{R}{\sqrt{2}}.
\]
(c) Substituindo $x=R/\sqrt2$:
\[
F_{\max}=\frac{2\,\ke Qq}{3\sqrt{3}\,R^{2}}
=\frac{2\cdot\pot{9}{9}\cdot\pot{1}{-5}\cdot\pot{2}{-6}}{3\sqrt3\cdot 0{,}01}
=\frac{0{,}36}{0{,}0520}=\SI{6.93}{\newton}.
\]
\textbf{Casos-limite:} em $x=0$ a força é nula (simetria perfeita); para
$x\gg R$, $F\to\ke Qq/x^{2}$ --- o anel se comporta como carga puntiforme.
Entre esses dois zeros deve existir um máximo, e ele está em $R/\sqrt2$.}
\end{questao}

\begin{questao}[4]{}
\textbf{(Distribuição contínua.)} Um bastão isolante retilíneo de comprimento
$L$ possui carga $Q$ uniformemente distribuída. Uma carga puntiforme $q$ está
sobre o prolongamento do bastão, a uma distância $d$ da extremidade mais
próxima.
\begin{itensq}
\item Deduza $F=\dfrac{\ke Qq}{d\,(d+L)}$.
\item Verifique o comportamento para $L\ll d$.
\item Calcule $F$ para $Q=\SI{4}{\micro\coulomb}$, $L=\SI{20}{\centi\meter}$,
      $q=\SI{1}{\micro\coulomb}$ e $d=\SI{10}{\centi\meter}$.
\end{itensq}
\resposta{(b) $F\to\ke Qq/d^{2}$; (c) $F=\SI{1.2}{\newton}$}
\resolucao{(a) Densidade linear $\lambda=Q/L$. Um elemento $\mathrm{d} x$ a uma
distância $x$ da carga carrega $\mathrm{d} Q=\lambda\,\mathrm{d} x$ e todas as contribuições
são \emph{colineares} --- por isso a integral é escalar:
\[
F=\int_{d}^{d+L}\frac{\ke q\lambda\,\mathrm{d} x}{x^{2}}
 =\ke q\lambda\left[-\frac{1}{x}\right]_{d}^{d+L}
 =\ke q\lambda\left(\frac{1}{d}-\frac{1}{d+L}\right)
 =\frac{\ke q\lambda L}{d(d+L)}=\frac{\ke Qq}{d(d+L)}.
\]
(b) Se $L\ll d$, então $d+L\approx d$ e $F\to\ke Qq/d^{2}$: recuperamos a lei de
Coulomb puntiforme, como deve ser.\\
(c) $F=\dfrac{\pot{9}{9}\cdot\pot{4}{-6}\cdot\pot{1}{-6}}{0{,}10\cdot 0{,}30}
=\dfrac{0{,}036}{0{,}03}=\SI{1.2}{\newton}$.\\
\textbf{Compare:} se toda a carga estivesse concentrada no centro do bastão
($\SI{20}{\centi\meter}$), teríamos $\SI{0.9}{\newton}$. A força real é maior
porque as porções mais próximas pesam com $1/x^{2}$ --- concentrar carga no
centro geométrico \emph{subestima} a força.}
\end{questao}

\begin{questao}[4]{}
\textbf{(Demonstração --- estabilidade.)} Duas cargas $+Q$ estão fixas em
$x=-d$ e $x=+d$. Uma carga $q$ é colocada na origem.
\begin{itensq}
\item Mostre que a origem é posição de equilíbrio para qualquer $q$.
\item Para $q>0$, analise a estabilidade a um pequeno deslocamento
      \emph{longitudinal} ($\pm x$).
\item Repita para um deslocamento \emph{transversal} ($\pm y$) e conclua.
\end{itensq}
\resposta{(a) por simetria; (b) estável, com $F\approx-4\ke Qq\,x/d^{3}$;
(c) instável, com $F\approx+2\ke Qq\,y/d^{3}$ --- não há equilíbrio estável em
todas as direções}
\resolucao{(a) As duas forças têm o mesmo módulo e sentidos opostos: resultante
nula, independentemente do sinal e do valor de $q$.\\
(b) Deslocando $q>0$ para $x$ pequeno:
\[
F_x=\ke Qq\left[\frac{1}{(d+x)^{2}}-\frac{1}{(d-x)^{2}}\right].
\]
Com $(d\pm x)^{-2}\approx d^{-2}(1\mp 2x/d)$:
\[
F_x\approx\frac{\ke Qq}{d^{2}}\left(-\frac{4x}{d}\right)=-\frac{4\ke Qq}{d^{3}}\,x.
\]
Sinal negativo $\Rightarrow$ força \textbf{restauradora}: estável nessa direção
(e o movimento é harmônico simples).\\
(c) Deslocando para $y$ pequeno, cada carga dista $\sqrt{d^{2}+y^{2}}\approx d$ e
repele $q$ com componente transversal $\ke Qq\,y/d^{3}$; as duas se somam:
\[
F_y\approx +\frac{2\ke Qq}{d^{3}}\,y,
\]
sinal positivo $\Rightarrow$ força que \textbf{afasta}: instável.\\
\textbf{Conclusão.} O equilíbrio é do tipo \emph{ponto de sela}: estável em uma
direção, instável na outra. Este é um caso particular do \textbf{teorema de
Earnshaw}: nenhuma configuração de cargas fixas pode manter uma carga puntiforme
em equilíbrio estável apenas por forças eletrostáticas. É por isso que modelos
puramente eletrostáticos do átomo não se sustentam.}
\end{questao}

\begin{questao}[4]{}
\textbf{(Série infinita.)} Cargas puntiformes de mesmo módulo $Q$ e
\textbf{sinais alternados} ocupam as posições $x=a,\,2a,\,3a,\dots$ sobre um
eixo: $-Q$ em $a$, $+Q$ em $2a$, $-Q$ em $3a$ e assim por diante,
indefinidamente. Uma carga $+q$ está na origem.
\begin{itensq}
\item Escreva a força resultante sobre $q$ como uma série.
\item Some a série e obtenha a expressão fechada.
\end{itensq}
\dados{$\displaystyle\sum_{n=1}^{\infty}\frac{(-1)^{n+1}}{n^{2}}=\frac{\pi^{2}}{12}$}
\resposta{$F=\dfrac{\pi^{2}}{12}\cdot\dfrac{\ke Qq}{a^{2}}
\approx 0{,}822\,\dfrac{\ke Qq}{a^{2}}$, atrativa (sentido $+x$)}
\resolucao{(a) A $n$-ésima carga está em $x=na$ e tem sinal $(-1)^{n}Q$. Uma
carga negativa \emph{atrai} $q$ (puxa no sentido $+x$); uma positiva
\emph{repele} (sentido $-x$). Tomando $+x$ como positivo:
\[
F=\frac{\ke Qq}{a^{2}}\left(\frac{1}{1^{2}}-\frac{1}{2^{2}}+\frac{1}{3^{2}}
-\frac{1}{4^{2}}+\cdots\right)
=\frac{\ke Qq}{a^{2}}\sum_{n=1}^{\infty}\frac{(-1)^{n+1}}{n^{2}}.
\]
(b) Usando o resultado dado:
\[
F=\frac{\pi^{2}}{12}\cdot\frac{\ke Qq}{a^{2}}\approx 0{,}822\,\frac{\ke Qq}{a^{2}}.
\]
\textbf{Interpretação.} Apesar de haver infinitas cargas, a força é
\emph{finita} e vale menos que a da primeira carga sozinha: a alternância de
sinais produz cancelamento progressivo, e o decaimento $1/n^{2}$ garante
convergência absoluta. Se todas as cargas fossem negativas, a soma seria
$\pi^2/6\approx1{,}645$ --- exatamente o dobro.}
\end{questao}

\begin{questao}[4]{}
\textbf{(Oscilações.)} Uma partícula de massa $m=\SI{1.0}{\gram}$ e carga
$-q=\SI{-2}{\micro\coulomb}$ é liberada em repouso sobre o eixo de um anel de
raio $R=\SI{10}{\centi\meter}$ e carga $Q=\SI{5}{\micro\coulomb}$, a uma
distância $x_0\ll R$ do centro.
\begin{itensq}
\item Mostre que, para $x\ll R$, o movimento é harmônico simples.
\item Obtenha $\omega$ e o período $T$.
\end{itensq}
\resposta{(b) $\omega=\sqrt{\ke Qq/(mR^{3})}=\SI{300}{\radian\per\second}$;
$T\approx\SI{20.9}{\milli\second}$}
\resolucao{(a) A força axial sobre a carga negativa é atrativa, dirigida ao
centro, de módulo $F=\ke Qq\,x/(x^{2}+R^{2})^{3/2}$. Para $x\ll R$, o
denominador tende a $R^{3}$:
\[
F_x\approx-\frac{\ke Qq}{R^{3}}\,x.
\]
Força proporcional ao deslocamento e de sinal oposto --- exatamente a forma
$F=-kx$ da lei de Hooke, com constante efetiva $k_{\text{ef}}=\ke Qq/R^{3}$.
Logo o MHS.\\
(b)
\[
\omega=\sqrt{\frac{k_{\text{ef}}}{m}}=\sqrt{\frac{\ke Qq}{mR^{3}}}
=\sqrt{\frac{\pot{9}{9}\cdot\pot{5}{-6}\cdot\pot{2}{-6}}
{\pot{1}{-3}\cdot(0{,}10)^{3}}}
=\sqrt{\frac{0{,}09}{\pot{1}{-6}}}=\SI{300}{\radian\per\second},
\]
\[
T=\frac{2\pi}{\omega}=\SI{20.9}{\milli\second}.
\]
\textbf{Observe:} o período não depende de $x_0$ --- é a assinatura do MHS. Fora
da aproximação $x\ll R$, a força deixa de ser linear e o movimento continua
periódico, mas com período dependente da amplitude.}
\end{questao}

\begin{questao}[4]{}
\textbf{(Modelagem com taxas.)} Duas esferas idênticas de massa
$m=\SI{2.0}{\gram}$, cada uma com carga $q$, pendem de fios isolantes de
comprimento $L=\SI{60}{\centi\meter}$ presos ao mesmo ponto. Elas se afastam até
uma separação $x\ll L$. Por fuga de carga para o ar úmido, as esferas se
aproximam lentamente.
\begin{itensq}
\item Mostre que $x^{3}=\dfrac{2\ke q^{2}L}{mg}$.
\item Demonstre que
      $\dfrac{\mathrm{d} x}{\mathrm{d} t}=\dfrac{2x}{3q}\dfrac{\mathrm{d} q}{\mathrm{d} t}$.
\item Se em certo instante $x=\SI{6.0}{\centi\meter}$ e as esferas se aproximam
      a $\SI{1.0}{\milli\meter\per\second}$, calcule $\mathrm{d} q/\mathrm{d} t$.
\end{itensq}
\dados{$g=\SI{9.8}{\meter\per\second\squared}$}
\resposta{(c) $q\approx\SI{19.8}{\nano\coulomb}$ e
$\mathrm{d} q/\mathrm{d} t\approx\SI{-0.49}{\nano\coulomb\per\second}$}
\resolucao{(a) Para cada esfera: tensão, peso e repulsão. Do triângulo de
forças, $\tan\theta=F/(mg)$. Com $x\ll L$,
$\tan\theta\approx\sin\theta=(x/2)/L$. Igualando:
\[
\frac{x}{2L}=\frac{\ke q^{2}}{x^{2}mg}
\;\Longrightarrow\;
x^{3}=\frac{2\ke q^{2}L}{mg}.
\]
(b) Derivando implicitamente em relação ao tempo:
\[
3x^{2}\frac{\mathrm{d} x}{\mathrm{d} t}=\frac{4\ke L}{mg}\,q\,\frac{\mathrm{d} q}{\mathrm{d} t}.
\]
De (a), $\dfrac{2\ke L}{mg}=\dfrac{x^{3}}{q^{2}}$; substituindo:
\[
3x^{2}\frac{\mathrm{d} x}{\mathrm{d} t}=\frac{2x^{3}}{q}\frac{\mathrm{d} q}{\mathrm{d} t}
\;\Longrightarrow\;
\frac{\mathrm{d} x}{\mathrm{d} t}=\frac{2x}{3q}\frac{\mathrm{d} q}{\mathrm{d} t}.
\]
(c) Primeiro $q$, por (a):
\[
q=\sqrt{\frac{x^{3}mg}{2\ke L}}
=\sqrt{\frac{(0{,}060)^{3}\cdot 0{,}0020\cdot 9{,}8}{2\cdot\pot{9}{9}\cdot0{,}60}}
=\pot{1.98}{-8}\,\si{\coulomb}.
\]
Com $\mathrm{d} x/\mathrm{d} t=-\pot{1.0}{-3}\,\si{\meter\per\second}$ (aproximação):
\[
\frac{\mathrm{d} q}{\mathrm{d} t}=\frac{3q}{2x}\frac{\mathrm{d} x}{\mathrm{d} t}
=\frac{3\cdot\pot{1.98}{-8}}{2\cdot 0{,}060}\cdot(-\pot{1}{-3})
=-\pot{4.9}{-10}\,\si{\coulomb\per\second}.
\]
\textbf{Leitura física:} medir uma \emph{velocidade} converte-se em medir uma
\emph{corrente de fuga} de meio nanoampère --- pequena demais para muitos
amperímetros, mas visível a olho nu pelo movimento das esferas.}
\end{questao}

\begin{questao}[4]{}
\textbf{(Simetria e integração.)} Um fio semicircular de raio
$R=\SI{5.0}{\centi\meter}$ possui carga $Q=\SI{3}{\micro\coulomb}$ distribuída
uniformemente. Uma carga $q=\SI{1}{\micro\coulomb}$ ocupa o centro do
semicírculo.
\begin{itensq}
\item Justifique por que a resultante é dirigida ao longo do eixo de simetria.
\item Mostre que $F=\dfrac{2\ke Qq}{\pi R^{2}}$ e calcule seu valor.
\end{itensq}
\resposta{$F=\SI{6.9}{\newton}$, ao longo do eixo de simetria}
\resolucao{(a) Para cada elemento do fio existe outro simétrico em relação ao
eixo vertical. Suas contribuições têm componentes perpendiculares ao eixo iguais
e opostas --- que se cancelam --- e componentes ao longo do eixo que se somam.\\
(b) Densidade linear $\lambda=Q/(\pi R)$. O elemento em $\theta$ tem
$\mathrm{d} Q=\lambda R\,\mathrm{d}\theta$, dista $R$ do centro e contribui com
$\ke q\,\mathrm{d} Q/R^{2}$; sua projeção no eixo vale $\sin\theta$:
\[
F=\int_{0}^{\pi}\frac{\ke q\lambda R\,\mathrm{d}\theta}{R^{2}}\sin\theta
 =\frac{\ke q\lambda}{R}\underbrace{\int_{0}^{\pi}\sin\theta\,\mathrm{d}\theta}_{=2}
 =\frac{2\ke q\lambda}{R}=\frac{2\ke Qq}{\pi R^{2}}.
\]
\[
F=\frac{2\cdot\pot{9}{9}\cdot\pot{3}{-6}\cdot\pot{1}{-6}}
{\pi\,(0{,}050)^{2}}=\frac{0{,}054}{\pot{7.85}{-3}}=\SI{6.9}{\newton}.
\]
\textbf{Compare:} se a mesma carga $Q$ estivesse concentrada em um ponto a $R$,
teríamos $\ke Qq/R^{2}=\SI{10.8}{\newton}$. O fator $2/\pi\approx 0{,}64$ mede
exatamente a perda por cancelamento vetorial ao longo do arco. Um anel
\emph{completo} levaria esse fator a zero.}
\end{questao}

\begin{questao}[4]{}
\textbf{(Síntese --- equilíbrio de um sistema.)} Quatro cargas iguais
$Q=\SI{2}{\micro\coulomb}$ ocupam os vértices de um quadrado de lado $a$. Uma
quinta carga $q$ é colocada no centro.
\begin{itensq}
\item Determine $q$ (sinal e módulo, em função de $Q$) para que \textbf{cada}
      carga dos vértices fique em equilíbrio.
\item Calcule $q$ numericamente.
\item O sistema completo está em equilíbrio estável? Justifique.
\end{itensq}
\resposta{(a) $q=-\dfrac{Q(2\sqrt{2}+1)}{4}$; (b) $q\approx\SI{-1.91}{\micro\coulomb}$;
(c) não --- o equilíbrio é instável}
\resolucao{(a) Por simetria, basta impor o equilíbrio de \emph{um} vértice; a
resultante das outras três cargas aponta ao longo da diagonal, para fora.
\begin{itemize}
\item Duas vizinhas, a distância $a$: cada uma contribui $\ke Q^{2}/a^{2}$, e
suas projeções sobre a diagonal valem $\cos\SI{45}{\degree}=\sqrt2/2$ cada, somando
$\sqrt{2}\,\ke Q^{2}/a^{2}$.
\item A oposta, a distância $a\sqrt{2}$: contribui
$\ke Q^{2}/(2a^{2})$, já ao longo da diagonal.
\end{itemize}
\[
F_{\text{vértices}}=\frac{\ke Q^{2}}{a^{2}}\left(\sqrt{2}+\frac{1}{2}\right).
\]
A carga central dista $a/\sqrt{2}$ do vértice, logo atua com
$F_{c}=\ke|q|Q/(a^{2}/2)=2\ke|q|Q/a^{2}$ e deve ser \textbf{atrativa} ---
portanto $q<0$. Igualando:
\[
2|q|=Q\left(\sqrt2+\frac12\right)
\;\Longrightarrow\;
|q|=\frac{Q\,(2\sqrt{2}+1)}{4}\approx 0{,}957\,Q.
\]
(b) $q\approx-0{,}957\cdot 2=\SI{-1.91}{\micro\coulomb}$.\\
(c) \textbf{Não.} A condição imposta zera a resultante sobre os vértices, mas
não sobre a carga central --- que já está em equilíbrio por simetria, embora
instável. Mais fundamentalmente, o teorema de Earnshaw proíbe equilíbrio
eletrostático estável: qualquer perturbação de $q$ ao longo de uma diagonal a
faz cair sobre um vértice. O equilíbrio existe, mas não sobrevive a uma
perturbação infinitesimal.}
\end{questao}
